\section{Trabajos futuros}
\label{sec:trabajos-futuros}

Los resultados obtenidos en la evaluación del prototipo, junto con las limitaciones descritas en la sección anterior, permiten identificar varias líneas de trabajo que se proyectan de forma natural a partir de este trabajo de grado y que resultan pertinentes para consolidar el aporte de FLP Viewer y ampliar su alcance pedagógico.

\begin{itemize}

\item \textbf{Evaluación con múltiples cohortes y otras instituciones.} La evaluación reportada en este documento se realizó con doce estudiantes de una única cohorte de la asignatura FLP, lo que restringe la posibilidad de generalizar los hallazgos de usabilidad y utilidad pedagógica. Una línea de trabajo futuro consiste en replicar la evaluación en cohortes posteriores y, de ser posible, en otros programas de Ingeniería de Sistemas que empleen Racket y la librería EOPL como sustrato de la asignatura, con el fin de contrastar los resultados obtenidos y fortalecer la evidencia empírica sobre el uso de este tipo de herramientas en contextos universitarios latinoamericanos, un vacío señalado en la revisión de antecedentes de este trabajo.

\item \textbf{Estudio longitudinal del impacto en el desempeño académico.} Los resultados favorables de usabilidad y utilidad pedagógica percibida reportados en el Capítulo~\ref{cap6} constituyen una base sólida para abordar una pregunta que este trabajo no cubrió, el efecto del prototipo sobre el desempeño académico y la retención del aprendizaje a lo largo del semestre. El diseño de un estudio longitudinal que compare el rendimiento de estudiantes que utilizan FLP Viewer frente a cohortes que no lo hacen permitiría medir con mayor rigor la contribución de la herramienta al aprendizaje de los contenidos de la asignatura.

\item \textbf{Retroalimentación enriquecida sobre errores de gramática.} En las preguntas abiertas del instrumento de evaluación, un grupo de estudiantes solicitó mensajes más detallados frente a errores de sintaxis en la definición de gramáticas BNF. El pipeline de gramática actual reporta la posición exacta del error, pero no sugiere correcciones ni explica la causa probable del fallo. Incorporar mecanismos de retroalimentación enriquecida, como la sugerencia de reglas de producción similares o la detección de errores comunes de la notación BNF, podría reducir la fricción reportada por este grupo de estudiantes.

\item \textbf{Extensión de las estrategias de visualización a la distinción entre ligadura y asignación.} La evaluación con usuarios identificó que el ítem relacionado con la distinción entre ligadura y asignación de variables concentró el menor nivel de acuerdo entre los siete evaluados, lo que confirma que esta diferencia conceptual, ya señalada como fuente de dificultad en el diagnóstico del Capítulo~\ref{cap3}, no queda completamente resuelta con la representación actual del panel de ambiente. Explorar formas complementarias de visualización, como el resaltado explícito de la celda modificada durante una operación \texttt{set} o una animación que distinga la creación de un nuevo frame de la actualización de una referencia existente, constituye una vía concreta para profundizar en este concepto.

\item \textbf{Persistencia y seguimiento del progreso del estudiante.} El prototipo actual no conserva información entre sesiones, ya que su diseño priorizó la interacción con la gramática, el AST y el ambiente de ejecución sobre la gestión de usuarios. Incorporar un mecanismo de identificación y almacenamiento del progreso permitiría que el docente responsable de la asignatura consulte el historial de interacción de cada estudiante con la biblioteca de ejemplos, información que resultaría valiosa tanto para el seguimiento pedagógico como para el estudio longitudinal propuesto en el punto anterior.

\item \textbf{Extensión a otros paradigmas de programación.} La encuesta aplicada a los docentes del programa evidenció dificultades similares en la enseñanza de programación lógica y de otros paradigmas avanzados, más allá del paradigma funcional abordado por FLP Viewer. Dado que el pipeline de gramática y las estrategias de visualización del prototipo no dependen de un paradigma específico, una línea de trabajo futuro consiste en evaluar su adaptación a intérpretes de lenguajes lógicos o imperativos, ampliando el alcance de la herramienta a otras asignaturas del programa que enfrentan dificultades conceptuales análogas.

\end{itemize}
